\chapter{终端的基本操作}
\label{ch:first-terminal}
\chapterlead{终端接受一行文本，交给程序执行，再把输出显示出来。本章说明提示符的构成、命令行的语法、编辑快捷键、报错的读法，以及环境变量中影响命令查找的 \texttt{PATH}。}

\section{起步：四条查询命令}

打开终端。Windows 可用 WSL，macOS 使用“终端”，多数 Linux 桌面的快捷键是 \texttt{Ctrl+Alt+T}。依次执行：

\begin{shellsession}
$ whoami
weiwen
$ pwd
/home/weiwen
$ date
2026 年 11 月 25 日 星期三 14 时 32 分 11 秒 CST
$ uptime
 14:32:11 up 3 days,  2:07,  1 user,  load average: 0.42, 0.31, 0.28
\end{shellsession}

四条命令分别给出：当前用户名、当前目录、当前时间、系统已运行时长与负载。负载的三个数字是 1、5、15 分钟的平均值，含义是等待处理的任务数；数值长期高于处理器核心数，说明处理器资源不足。

这些命令由 shell 执行。shell 读取输入的一行，找到对应的程序，传递参数，再把结果写回终端。

\section{提示符的构成}

交互式 shell 在等待输入时显示提示符，通常形如：

\begin{center}
\figfit{%
\begin{tikzpicture}[font=\ttfamily\small]
  \node[anchor=west] (u) at (0,0) {\textcolor{teal}{weiwen}};
  \node[anchor=west] (at) at ([xshift=1mm]u.east) {\textcolor{muted}{@}};
  \node[anchor=west] (h) at ([xshift=1mm]at.east) {\textcolor{teal}{nixos}};
  \node[anchor=west] (c) at ([xshift=1mm]h.east) {\textcolor{muted}{:}};
  \node[anchor=west] (p) at ([xshift=1mm]c.east) {\textcolor{ink}{\textasciitilde/projects}};
  \node[anchor=west] (d) at ([xshift=3mm]p.east) {\textcolor{muted}{\$}};
  \node[font=\scriptsize, text=muted, below=6mm of u] {用户名};
  \node[font=\scriptsize, text=muted, below=6mm of h] {主机名};
  \node[font=\scriptsize, text=muted, below=6mm of p] {当前目录};
  \node[font=\scriptsize, text=muted, below=6mm of d] {权限标记};
  \draw[rule] (u.south) -- ++(0,-4mm);
  \draw[rule] (h.south) -- ++(0,-4mm);
  \draw[rule] (p.south) -- ++(0,-4mm);
  \draw[rule] (d.south) -- ++(0,-4mm);
\end{tikzpicture}
}
\end{center}

三点需要记住。\texttt{\textasciitilde} 表示家目录，\texttt{\textasciitilde/projects} 即“家目录下的 projects”。末尾的 \texttt{\$} 表示普通用户，\texttt{\#} 表示 root。提示符不是命令的一部分，抄写命令时不应包含它。

若提示符变为 \texttt{>}，表示上一条命令缺少结尾，通常由引号未配对引起。按 \texttt{Ctrl+C} 放弃该行并重新输入。

\section{命令的结构}

\begin{shellsession}
$ ls -lh --sort=size /var/log
\end{shellsession}

这一行由四部分组成。

\begin{center}
\figfit{%
\begin{tikzpicture}[font=\ttfamily\small, node distance=1mm]
  \node[draw=teal!70, fill=teal!7, rounded corners=1pt, inner sep=3pt] (cmd) {ls};
  \node[draw=ochre!70, fill=ochre!7, rounded corners=1pt, inner sep=3pt, right=3mm of cmd] (opt) {-lh};
  \node[draw=ochre!70, fill=ochre!7, rounded corners=1pt, inner sep=3pt, right=1mm of opt] (opt2) {--sort=size};
  \node[draw=ink!40, fill=ink!5, rounded corners=1pt, inner sep=3pt, right=3mm of opt2] (arg) {/var/log};
  \node[font=\scriptsize, text=teal, below=3mm of cmd] {命令};
  \node[font=\scriptsize, text=ochre, below=3mm of opt2] {选项};
  \node[font=\scriptsize, text=ink, below=3mm of arg] {参数};
\end{tikzpicture}
}
\end{center}

\begin{itemize}
  \item \textbf{命令}：要运行的程序名。shell 在 \texttt{PATH} 列出的目录中查找它。
  \item \textbf{选项}：改变程序行为。单横线为短选项，可以合并（\texttt{-l -h} 写作 \texttt{-lh}）；双横线为长选项，通常用等号赋值。
  \item \textbf{参数}：命令作用的对象，上例中为目录 \texttt{/var/log}。
\end{itemize}

空格是分隔符。\texttt{ls -lh} 与 \texttt{ls-lh} 含义不同，后者会被当作一个不存在的命令名。路径中含空格时需加引号。

\begin{shellsession}
$ cd "My Documents"      # 正确
$ cd My Documents        # 被解析为两个参数：My 与 Documents
\end{shellsession}

\section{补全与历史}

\begin{itemize}
  \item \textbf{Tab}：补全文件名与命令名。输入 \texttt{cd /var/lo} 后按 Tab，补全为 \texttt{/var/log}；连按两次 Tab 列出所有候选。
  \item \textbf{上箭头}：调出上一条命令，可连续向前翻阅。
  \item \textbf{Ctrl+R}：反向搜索历史。输入 \texttt{ssh} 后显示最近包含该串的命令，再次按 \texttt{Ctrl+R} 向前搜索。
  \item \textbf{!!}：重复上一条命令。\texttt{sudo !!} 表示以 root 权限重跑上一条命令。
  \item \textbf{!\_\$}：引用上一条命令的最后一个参数。\texttt{mkdir demo} 之后执行 \texttt{cd !\_\$} 可直接进入该目录。
\end{itemize}

补全与历史的作用不止于节省时间：它们同时避免了拼写错误引起的误操作。

\section{行编辑：readline}

提示符下输入的每一行，由行编辑库处理。bash 用的是 GNU Readline；zsh 用自己的 ZLE，默认同样是 Emacs 键位。Python 的交互解释器、\texttt{psql}、\texttt{mysql} 客户端也使用 readline 或等价实现。这套键位大部分与 Emacs 相同，两个例外在下面的表后说明。

\begin{booktable}{\textwidth}{L{0.72} L{0.94} J{1.34}}
\toprule
键位 & 动作 & 用途 \\
\midrule
\endhead
\texttt{Ctrl+a} & 移到行首 & 改一条长命令的开头 \\
\texttt{Ctrl+e} & 移到行尾 & 回到末尾继续输入 \\
\texttt{Ctrl+b}、\texttt{Ctrl+f} & 左移、右移一个字符 & 逐字调整 \\
\texttt{Alt+b}、\texttt{Alt+f} & 左移、右移一个词 & 按词跳转，比逐字快 \\
\texttt{Ctrl+k} & 删除光标到行尾 & 丢掉后半段 \\
\texttt{Ctrl+u} & 删除光标到行首 & 丢掉前半段 \\
\texttt{Ctrl+w} & 删除前一个词 & 相当于删一个参数 \\
\texttt{Alt+d} & 删除后一个词 & 局部修改 \\
\texttt{Ctrl+y} & 粘贴最近删除的内容 & 与上四行配合移动文本 \\
\texttt{Ctrl+t} & 交换两个字符 & 修正手误，不必重输 \\
\texttt{Ctrl+l} & 清屏 & 输出太乱时 \\
\bottomrule
\end{booktable}

两个键位与前文的 \texttt{Ctrl+c}、\texttt{Ctrl+r} 同属这套体系：\texttt{Ctrl+c} 放弃整行，\texttt{Ctrl+r} 反向搜索历史。删除与粘贴配合可以搬运文本：\texttt{Ctrl+u} 删掉光标前半段，按 \texttt{Ctrl+e} 移到行尾，再按 \texttt{Ctrl+y} 贴回。

表中有两个键在 Emacs 里含义不同：\texttt{Ctrl+u} 是前缀参数，\texttt{Ctrl+w} 是剪切选区。其余键位两边一致。

\begin{notebox}[Ctrl+s 被终端占用]
readline 里 \texttt{Ctrl+s} 是向前搜索历史，与 \texttt{Ctrl+r} 相对。终端驱动的流控（\texttt{IXON}）会截住 \texttt{Ctrl+s} 用于暂停输出，按下后屏幕不再刷新，看起来像卡死。按 \texttt{Ctrl+q} 恢复。要长期启用，把 \texttt{stty -ixon} 写进 shell 的启动文件。
\end{notebox}

\subsection{用编辑器写长命令}

一行命令写到一半，结构复杂到不便在提示符里改时，按 \texttt{Ctrl+x} \texttt{Ctrl+e}。readline 会把当前行交给 \texttt{\$EDITOR} 指定的编辑器，保存退出后命令立即执行。

\begin{shellsession}
$ echo "${VISUAL:-（未设置）}" "${EDITOR:-（未设置）}"
（未设置） （未设置）
$ export EDITOR=emacs
\end{shellsession}

readline 先取 \texttt{VISUAL}，再取 \texttt{EDITOR}；两者都未设置时用 \texttt{emacs}。变量必须导出，仅赋值不会被 readline 读到。设为 \texttt{emacs} 之后，\texttt{Ctrl+x} \texttt{Ctrl+e} 会打开 Emacs；终端里启动 Emacs 的方法在第\ref{ch:editor-workflow}章说明。

\subsection{inputrc}

readline 的行为由 \texttt{\textasciitilde/.inputrc} 调整，不必改 shell 脚本。

\begin{textcode}
set editing-mode emacs
set completion-ignore-case on
set bell-style none

$if Bash
  "\C-x\C-r": re-read-init-file
$endif
\end{textcode}

三行常规设置：保持 Emacs 键位、补全忽略大小写、关闭提示音。最后一段只在 bash 中生效，把 \texttt{Ctrl+x} \texttt{Ctrl+r} 绑成重新读取本文件，改完不必退出 shell。

这份文件同样可以交给 Home Manager 生成，写法见第\ref{ch:editor-workflow}章。

\section{中断、挂起与退出}

三个按键的含义不同。

\begin{booktable}{\textwidth}{L{1.114} L{1.029} J{0.858}}
\toprule
按键 & 动作 & 使用场景 \\
\midrule
\endhead
\texttt{Ctrl+C} & 中断当前程序 & 命令无响应、输出过多、参数写错 \\
\texttt{Ctrl+Z} & 挂起到后台 & 临时返回提示符，稍后继续 \\
\texttt{Ctrl+D} & 发送输入结束 & 退出 shell，或结束等待输入的程序 \\
\bottomrule
\end{booktable}

挂起后用 \texttt{jobs} 查看后台任务，\texttt{fg} 将最近的任务调回前台，\texttt{bg} 使其在后台继续运行。三者之中最常用的是 \texttt{Ctrl+C}。

\begin{warnbox}[Ctrl+C 不能撤销已完成的写入]
它只中断仍在前台运行的程序。已经写入磁盘的数据不会恢复。删除文件、格式化磁盘这类操作，中断的时机通常已经太晚；执行前应确认路径与参数。
\end{warnbox}

\section{报错的四类形态}

报错信息包含定位线索，常见形态有四种。

\begin{shellsession}
$ 1s
bash: 1s: command not found
$ cat /etc/shadow
cat: /etc/shadow: Permission denied
$ cat notes.txt
cat: notes.txt: No such file or directory
$ echo "hello
>
\end{shellsession}

\begin{itemize}
  \item \textbf{command not found}：未找到该程序。常见原因是拼写错误（上例将字母 \texttt{l} 写成了数字 \texttt{1}），或程序未安装、不在 \texttt{PATH} 中。
  \item \textbf{Permission denied}：程序存在，但权限不足。用 \texttt{ls -l} 查看文件权限，第\ref{ch:permissions}章详述。
  \item \textbf{No such file or directory}：路径错误。用 \texttt{pwd} 确认当前目录，用 \texttt{ls} 确认文件是否存在。
  \item 提示符变为 \texttt{>}：引号或括号未闭合。补齐后继续，或按 \texttt{Ctrl+C} 放弃该行。
\end{itemize}

读报错可按固定顺序：先判断类别，再查看涉及的路径或名称，最后才检索。直接复制整段报错搜索，通常得到多个互不相关的答案。

\section{PATH 与环境变量}

shell 中有一组变量影响程序行为，最先需要掌握的是 \texttt{PATH}。

\begin{shellsession}
$ echo $PATH
/home/weiwen/.nix-profile/bin:/run/current-system/sw/bin:/usr/bin:/bin
$ command -v rg
/home/weiwen/.nix-profile/bin/rg
\end{shellsession}

\texttt{PATH} 是用冒号分隔的目录列表。输入 \texttt{rg} 时，shell 按顺序在这些目录中查找同名文件，执行第一个匹配项。程序装在 \texttt{PATH} 之外时同样会出现该错误。因此“命令找不到”有两个常见原因：程序未安装，或不在 \texttt{PATH} 中。

\begin{shellsession}
$ export EDITOR=nvim      # 仅对当前 shell 有效
$ echo $EDITOR
nvim
\end{shellsession}

\texttt{export} 设置的变量会传递给随后启动的程序。需要长期生效时，写入 shell 的配置文件；在 NixOS 上推荐通过 Home Manager 声明（第\ref{ch:flakes-homemanager}章）。

\section{shell 的选择}

常见 shell 有三种。

\begin{booktable}{\textwidth}{L{1.064} L{1.064} J{0.871}}
\toprule
shell & 特点 & 适用场景 \\
\midrule
\endhead
bash & 预装于多数系统，语法最通用 & 编写脚本、登录服务器、照抄教程 \\
zsh & 补全能力强，主题丰富 & 日常交互，愿意维护配置 \\
fish & 默认开启语法高亮与自动建议 & 不打算配置交互环境 \\
\bottomrule
\end{booktable}

本书的命令在三种 shell 下均可执行。脚本一律用 bash 编写：服务器上必然存在 bash，不一定存在 zsh 或 fish。交互环境可按偏好选择，脚本选择最通用的实现。

\section{旧工具与新工具}

\begin{booktable}{\textwidth}{L{1.190} L{0.967} J{0.843}}
\toprule
旧工具 & 新工具 & 差异 \\
\midrule
\endhead
\texttt{cd} & \texttt{zoxide} & 记录访问过的目录，\texttt{z proj} 直接跳转 \\
\texttt{ls} & \texttt{eza} & 颜色、图标、树状视图、Git 状态 \\
\texttt{cat} & \texttt{bat} & 语法高亮、行号、自动分页 \\
\texttt{history} & \texttt{atuin} & 跨机器同步，可按目录与退出码过滤 \\
\bottomrule
\end{booktable}

\begin{lab}[练习：四个按键]
\begin{enumerate}
  \item 输入 \texttt{cd /usr/} 后按 Tab，记录补全结果；再按一次 Tab，记录列出的候选。
  \item 运行 \texttt{sleep 30}，用 \texttt{Ctrl+C} 中断。
  \item 运行 \texttt{sleep 30}，用 \texttt{Ctrl+Z} 挂起，\texttt{jobs} 查看，\texttt{fg} 调回，再用 \texttt{Ctrl+C} 结束。
  \item 用 \texttt{Ctrl+R} 找出十分钟前输入的第一条命令。
\end{enumerate}
\end{lab}

\section{练习}

\begin{enumerate}
  \item 用 \texttt{echo} 输出家目录路径。可用变量 \texttt{\$HOME}。
  \item 统计 \texttt{PATH} 中包含几个目录。可将 \texttt{echo \$PATH} 的输出交给 \texttt{tr ':' '\textbackslash n'} 后计数。
  \item 制造三类报错：命令名拼错、访问无权限文件、打开不存在的文件，把三条报错原文抄入笔记。
  \item 分别用 \texttt{type} 与 \texttt{command -v} 查询同一命令，比较两者输出的差异。
  \item 输入一条长命令，用 \texttt{Ctrl+a}、\texttt{Ctrl+k}、\texttt{Ctrl+y} 把它拆成两半再拼回去。
\end{enumerate}

\begin{checkpoint}
本章结束时，应能：
\begin{itemize}
  \item 说出提示符中用户名、主机名、当前目录与权限标记的位置；
  \item 将一条命令拆分为命令、选项、参数三部分；
  \item 使用 Tab、命令历史与 \texttt{Ctrl+R} 减少输入量；
  \item 用 readline 键位在行内移动、删除与粘贴，知道这套键位与 Emacs 同源；
  \item 用 \texttt{Ctrl+C}、\texttt{Ctrl+Z}、\texttt{Ctrl+D} 控制前台程序；
  \item 判断四类常见报错的原因；
  \item 解释 \texttt{PATH} 的作用，说明“命令找不到”的常见原因。
\end{itemize}
\end{checkpoint}
